\section{Method}
\label{sec:method}

\subsection{Overview of the research design}
\label{sec:design}

The study is designed as empirical research that verifies medical VQA domain
adaptation of lightweight VLMs sequentially through a three-phase experiment.
Phase~1 measures the zero-shot performance of the four models before fine-tuning
in order to answer RQ1 (differences across models). Phase~2 applies QLoRA
fine-tuning to verify RQ2 (the effect of fine-tuning) while simultaneously
exploring the optimal QLoRA configuration (data scale, LoRA rank, target modules)
through ablation studies. Phase~3 takes the optimal configuration identified in
Phase~2 as its starting point and compares four HPO strategies --- Manual, Random
Search, Optuna (TPE), and Autoresearch --- to verify RQ3 (the competitiveness of
autonomous search). Each phase has a sequential structure in which the results of
the preceding phase (best model, optimal QLoRA configuration) become fixed
conditions for the next.

\subsection{Experimental environment}
\label{sec:hardware}

All experiments reported here were conducted on a cloud GPU service (RunPod).
Phases~1 and 2 used an RTX 4090 (24\,GB VRAM) instance and Phase~3 an RTX 3090
(24\,GB VRAM) instance. The change of GPU in Phase~3 was not a design choice but
the result of securing whatever instance was available under tightened budget
constraints. Both are consumer-class cards with 24\,GB of VRAM, and every Phase~3
trial ran on the same 3090 instance, so the hardware conditions of the strategy
comparison were controlled. A local workstation (RTX 5060 Ti 16\,GB $+$ RTX 4060
8\,GB, Ryzen 5 5600X, 32\,GB RAM) was used only for a preliminary smoke test to
size the Phase~3 execution; none of the figures reported in this paper were
produced in the local environment.

Feasibility in the 16\,GB-class environment targeted by this work was confirmed
\emph{from measured VRAM usage rather than by execution on such a card}. Peak
VRAM during Phase~2 QLoRA training was at most 14,373\,MB across the four models
(Gemma4-E2B), within the 16\,GB limit, while the other three models required only
4,015--7,943\,MB.

Model loading and QLoRA fine-tuning used the HuggingFace \texttt{transformers}
library together with the \texttt{unsloth} backend for accelerated 4-bit
quantized training, with LoRA adapters configured through the \texttt{peft}
library. The Bayesian optimization control in Phase~3 used \texttt{Optuna} (TPE
sampler), and experiment management and logging used \texttt{wandb}. BERTScore
computation for scoring open-ended responses used the \texttt{bert-score} library
(roberta-large / BioBERT backbones), and statistical analysis used \texttt{scipy}
and \texttt{statsmodels} (mixed-effects model). The LLM agent of the Autoresearch
strategy calls the Anthropic Claude API.

\subsection{Target models and datasets}
\label{sec:models}

The four models examined and the grounds for their selection are given in
Table~\ref{tab:models}. The selection criteria were threefold: (1) QLoRA
fine-tuning must be feasible within 16\,GB of VRAM (an \emph{a priori}
expectation, confirmed \emph{post hoc} from measured peak VRAM,
Section~\ref{sec:hardware}); (2) the license must be Apache 2.0 or MIT so that
research use is unrestricted; and (3) the model must have sufficient community and
framework support. Gemma4-E2B was included because its mixture-of-experts-family
PLE technique activates only 2.3B parameters at inference while providing
5.1B-class expressive capacity.

\begin{table}[t]
\centering
\caption{Target models.}
\label{tab:models}
\small
\begin{tabular}{@{}llll@{}}
\toprule
Model & Parameters & Architectural features & Expected QLoRA VRAM \\
\midrule
Qwen3-VL-2B    & 2B   & Thinking mode, DeepStack             & $\sim$8--10\,GB \\
Qwen2.5-VL-3B  & 3B   & Dynamic resolution, OCR in 19 languages & $\sim$8--10\,GB \\
SmolVLM2-2.2B  & 2.2B & HuggingFace lightweight VLM          & $\sim$8--10\,GB \\
Gemma4-E2B     & 2.3B (active) / 5.1B (total) & PLE (per-layer embeddings), Apache 2.0 & $\sim$12--14\,GB \\
\bottomrule
\end{tabular}
\end{table}

The three public medical VQA datasets are summarized in
Table~\ref{tab:datasets}. Each dataset was used with its official
train/val/test split unchanged, and all reported accuracies are computed on the
\emph{test} split only. Because the size of a distributed copy can differ from
the figure reported in the dataset's original paper, Table~\ref{tab:datasets}
separates the three quantities that are easy to conflate: the size reported by
the source publication, the size of the copy we actually obtained, and the
number of items each model was actually evaluated on.

\begin{table}[t]
\centering
\caption{Target datasets. ``Published'' is the QA-pair count reported by the
dataset's own paper; ``Obtained'' is the total size of the distributed copy used
here; ``Evaluated'' is the test-split size on which every reported accuracy is
computed. No subsampling was applied --- the evaluated counts are the full test
splits.}
\label{tab:datasets}
\small
\begin{tabular}{@{}llrrrl@{}}
\toprule
Dataset & Source distribution & Images & Published & Obtained & Evaluated \\
\midrule
PathVQA & \texttt{flaviagiammarino/path-vqa}      & 4,998 & 32,799 & 32,632 & 6,719 \\
SLAKE   & \texttt{mdwiratathya/SLAKE-vqa-english} & 642   & 14,028 & 7,033  & 1,061 \\
VQA-RAD & \texttt{flaviagiammarino/vqa-rad}       & 315   & 2,248  & 2,244  & 451 \\
\bottomrule
\end{tabular}
\end{table}

PathVQA covers pathology tissue images with seven question types; SLAKE and
VQA-RAD cover radiology and CT/MRI/X-ray images. All three mix open-ended and
closed-ended questions. Note that the SLAKE distribution used here is the
English-only release: the 14,028 figure published for SLAKE counts its bilingual
English--Chinese pairs, whereas the copy obtained for these experiments contains
the English half alone.

\paragraph{Contamination control.}
PathVQA (2018), SLAKE (2021), and VQA-RAD (2018) were all released before the
pretraining cut-off of the target models (2025--2026), so the possibility of
pretraining data contamination cannot be excluded. We measure such contamination
actively using the Min-$K$\% Probability attack \citep{shi2024minkprob}. The mean
token-level log-probability of the lowest $K$\% ($K = 20$) of tokens in each
sample's gold answer text is used as the contamination indicator --- on the
theory that samples exposed during pretraining exhibit higher mean probability
--- and samples in the top 5\% within a dataset are classified as suspected
contamination, after which the principal conclusion (RQ1) is re-verified on the
reduced sample set with those samples removed. The interpretation criteria are
that a difference of less than 1\,\%p between original and reduced results is
deemed robust, 1--5\,\%p is stated as a limitation, and more than 5\,\%p requires
reconsideration of the conclusion.

\subsection{Phase 1: zero-shot baseline evaluation}
\label{sec:phase1}

Zero-shot performance before fine-tuning was measured for $4$ models $\times$ $3$
datasets $=$ 12 conditions. Because evaluation uses greedy decoding it is
deterministic: changing the seed does not change the result, so repeated trials
are meaningless. Evaluation therefore used a single seed (42), and uncertainty is
reported as a bootstrap 95\% confidence interval over the per-sample
correct/incorrect decisions of each condition. Since the four models are
evaluated on an identical test set, comparison across models has a paired
structure; accordingly, Cochran's $Q$ test (binary correctness on a shared test
set, $H_0$: equal accuracy) and pairwise McNemar post-hoc tests with Bonferroni
correction were used instead of ANOVA, which assumes independent samples. The
measured metrics are closed-ended accuracy (multiple choice), open-ended accuracy
(gold-token matching) together with BERTScore F1 (roberta-large), bootstrap 95\%
CIs for each accuracy, response time (ms per item), and peak VRAM (MB).

\subsection{Phase 2: QLoRA fine-tuning}
\label{sec:phase2}

The base QLoRA configuration applied to every Phase~2 condition is given in
Table~\ref{tab:qlora}. The training budget targeted 3 epochs, but a ceiling of
\texttt{max\_steps}~$=500$ (samples\_seen fixed at 4,000 per condition) was
imposed because of cloud GPU time and cost constraints. Since the training volume
is fixed regardless of dataset size, the small VQA-RAD receives roughly 2 epochs
or more while the medium SLAKE and large PathVQA receive less than 1 epoch; this
asymmetry is a stated limitation of the study and is carried into the
interpretation of the results.

\begin{table}[t]
\centering
\caption{Base QLoRA configuration (Phase 2).}
\label{tab:qlora}
\small
\begin{tabular}{@{}ll@{}}
\toprule
Parameter & Value \\
\midrule
Quantization    & NF4 (4-bit NormalFloat) \\
LoRA rank       & 64 (determined by Ablation B) \\
LoRA alpha      & 128 \\
LoRA dropout    & 0.05 \\
Target modules  & all-linear (determined by Ablation C) \\
Learning rate   & 2e-4 \\
Batch size      & 1 (gradient accumulation 8, effective batch $=$ 8) \\
Optimizer       & paged\_adamw\_8bit \\
\bottomrule
\end{tabular}
\end{table}

The experimental conditions are $4$ models $\times$ $3$ datasets $=$ 12
conditions, each repeated 3 times (seeds 42/123/456). The values LoRA rank $=64$
and target $=$ all-linear were fixed on the basis of three ablation studies, each
with PathVQA and Qwen3-VL-2B held fixed: Ablation~A (effect of data size) varies
the training-data ratio over 5/10/25/50/100\%; Ablation~B (effect of LoRA rank)
varies rank $\in \{4, 8, 16, 32, 64\}$; and Ablation~C (effect of target modules)
compares $\{$q/v\_proj$\}$ versus $\{$q/k/v/o\_proj$\}$ versus
$\{$all-linear$\}$.

Catastrophic forgetting is measured in two ways. (A) Change in general
capability: the rate of accuracy decrease before and after fine-tuning is
measured on a VQAv2 validation subset (2,000 samples) across all 12 conditions.
(B) Cross-dataset generalization within the medical domain: evaluation on a
dataset other than the training dataset (for example, train on PathVQA $\to$
evaluate on SLAKE/VQA-RAD), giving $12$ conditions $\times$ $2$ cross-datasets
$=$ 24 additional evaluations. Because PathVQA (pathology) and SLAKE/VQA-RAD
(radiology) differ in image domain itself, the results of (B) are interpreted as
a domain-generalization gap rather than as catastrophic forgetting in the strict
sense.

\subsection{Phase 3: autonomous hyperparameter optimization}
\label{sec:phase3}

The hyperparameter search space shared by the four strategies is given in
Table~\ref{tab:searchspace}. The four strategies compared are Manual (the
researcher's default values, one run), Random Search (random sampling), Optuna
TPE (Bayesian optimization), and Autoresearch (autonomous search by an LLM
agent). Autoresearch repeats a loop that (1) reads the previous experimental
results (\texttt{results.tsv}), (2) proposes the next configuration together with
a natural-language rationale (\texttt{config.yaml} $+$ \texttt{rationale.md}),
(3) performs a git commit, (4) runs fixed training, (5) evaluates on the
validation set, and (6) retains the configuration if performance improves and
discards it otherwise.

\begin{table}[t]
\centering
\caption{Phase 3 hyperparameter search space.}
\label{tab:searchspace}
\small
\begin{tabular}{@{}lll@{}}
\toprule
Parameter & Search range & Type \\
\midrule
lora\_rank        & $\{4, 8, 16, 32, 64\}$        & Discrete \\
lora\_alpha       & rank $\times \{1, 2, 4\}$     & Discrete \\
learning\_rate    & $[1\mathrm{e}{-5},\ 5\mathrm{e}{-4}]$ & Continuous (log scale) \\
batch\_size       & $\{1, 2, 4\}$                 & Discrete \\
grad\_accum\_steps & $\{4, 8, 16\}$               & Discrete \\
warmup\_ratio     & $[0.0,\ 0.1]$                 & Continuous \\
weight\_decay     & $[0.0,\ 0.1]$                 & Continuous \\
lora\_targets     & $\{$minimal, medium, full$\}$ & Categorical \\
\bottomrule
\end{tabular}
\end{table}

All trials share the same model (the best model from Phase~2), the same dataset
(PathVQA), and a fixed \texttt{max\_steps}~$=200$ in order to control training
volume; the wall-clock ceiling \texttt{time\_budget\_min} serves as a safety
device to prevent pathological combinations rather than as an experimental
control variable. The original design called for Manual 10 $+$ Random Search 400
$+$ Optuna 400 $+$ Autoresearch 400 $=$ 1,210 trials in total (40 trials per
strategy $\times$ 10 independent repeats). A preliminary smoke test conducted
solely to size the execution indicated on a wall-clock basis --- counting not
only training but also validation and final test evaluation time --- that the
original scale would require roughly 24--25 days even with two GPUs in parallel.
The number of repeats (10), which is the unit of statistical testing and the
basis of run-level statistical power, was therefore preserved, while the number
of search trials per strategy was reduced from 40 to 20, halving the total time
to approximately 12.8 days. The execution scale of the four-strategy comparison
is thus Manual 10 $+$ Random Search 200 $+$ Optuna 200 $+$ Autoresearch 200 $=$
610 trials in total (20 trials per strategy $\times$ 10 independent repeats).
This reduction carries the trade-off of halving the diversity of hyperparameter
combinations explored per strategy, but it does not affect the validity of the
run-level statistical tests based on 10 independent repeats.

In addition, a 200-trial configuration-consistency re-experiment
(Autoresearch-v2) was conducted separately. This is not part of the four-strategy
same-condition comparison but a follow-up intervention designed to test whether
the duplicate-proposal phenomenon dissolves once the prompt--implementation
inconsistencies are removed. The search space, model, dataset,
\texttt{max\_steps}, and repetition count were held identical to the original,
and only the five configuration inconsistencies were corrected; the original
condition was preserved intact for reproducibility. The total execution scale of
Phase~3 is 810 trials.

Statistical verification is performed only at the run level, not at the trial
level. Because Autoresearch and Optuna are sequential optimizers, trials within
the same run are dependent (the result of trial $t$ influences the proposal at
$t+1$), violating the assumption of independent observations. The unit of testing
is the 10 final performance values obtained from the 10 independent repeats of
each strategy, analysed with the Kruskal--Wallis test (four-group comparison),
the Mann--Whitney $U$ test (pairwise, Autoresearch vs.\ Optuna), and BCa
bootstrap 95\% CIs. Trial-level data are used only for visualization such as
anytime performance curves.

\subsection{Evaluation metrics and statistical analysis}
\label{sec:stats}

\paragraph{Dual reporting of BERTScore.}
Open-ended responses are reported with both exact match and BERTScore F1. The
general-purpose criterion (roberta-large, threshold $\geq 0.7$) is the sole
decision metric (primary) for accuracy and statistical testing, while the
medical-specialized criterion (BioBERT, \texttt{dmis-lab/biobert-v1.1}) is
reported alongside as a secondary metric only; no dual gating (requiring both
metrics to pass before an answer is judged correct) is applied.

\paragraph{Dual measurement of catastrophic forgetting.}
Both (A) the change in general capability measured on VQAv2 and (B) the
cross-dataset generalization gap described in Section~\ref{sec:phase2} are
reported, so that side effects of fine-tuning that a single metric cannot capture
are measured from multiple angles.

\paragraph{Clinical significance analysis (WCA $+$ ECE).}
Motivated by the concern that plain accuracy does not fully capture the clinical
value of medical AI, weighted clinical accuracy (WCA) is computed using the
labels of the seven PathVQA question types (diagnosis, location, measurement,
description, temporal, yes\_no, unknown):
\[
\mathrm{WCA} = \frac{\sum (\text{accuracy per type} \times \text{weight})}{\sum \text{weights}} .
\]
Weights were assigned by clinical importance in the order diagnosis $=1.0$ (a
diagnostic error directly affects the direction of treatment) $>$ location $=0.8$
$>$ measurement $=0.7$ $>$ description $=0.6$ $>$ temporal $=$ yes\_no $=0.5$
(binary judgements with limited information content). These weights were set
arbitrarily by the researcher without external clinical literature or Delphi
consensus; they cannot be interpreted as an absolute scale of clinical importance
and are used only in a limited way as a supplementary reference metric
complementing the primary metrics (accuracy and BERTScore). Expected calibration
error (ECE) \citep{guo2017calibration}, a metric measuring how well a model's
predicted confidence matches its actual accuracy, was planned for inclusion but
could not be computed because the evaluation pipeline does not store per-sample
confidence.

\paragraph{Robust statistics (bootstrap $+$ mixed effects).}
Given the sample-size limitation of $n = 9$ (3 datasets $\times$ 3 seeds), the
test of the fine-tuning effect in Phase~2 uses three methods in parallel: (1) a
paired $t$-test with Cohen's $d$ (the conventional comparison), (2) a 95\% CI for
Cohen's $d$ obtained by BCa bootstrap with 10,000 resamples (robust estimation),
and (3) the Wilcoxon signed-rank test (non-parametric). A mixed-effects model
pooling the four models (\texttt{accuracy $\sim$ condition $+$ dataset}, group
$=$ seed) is additionally applied as a supplement; however, because heterogeneous
effects across models may cancel in a pooled estimate, the triple verification
performed per model is treated as the primary evidence and the mixed-effects
result is used only as supporting explanation --- a phenomenon actually observed
in Section~4.2.1.
